\section{模型评价与改进}

本文的评价对象是由数据结构、模型边界和求解规则共同构成的整体框架，而非孤立评价某一算法。针对稀疏、非正交且缺少重复测量的实验数据，框架优先保证结论可追溯，再在相应边界内追求信息利用效率。

\subsection{模型优点}

\begin{enumerate}
  \item \textbf{本模型将数据层级与问题目标一一对应。} 问题一只在同一催化剂组合内刻画温度响应，问题二才在严格匹配背景下比较因素，问题三再把转化率和选择性合成为收率优化目标。由此避免将 114 条温度记录误作 114 个独立催化剂设计，也避免以单项响应替代收率判断。
  \item \textbf{本模型以低复杂度规则约束小样本拟合。} 问题一的一次--二次候选、组内留一预测和物理边界检查相互制衡；问题三用分段线性连续化保持相邻节点连续而不制造内部峰值。该设计把“拟合得更好”与“可用于排序”明确区分，符合小样本交叉验证应审慎解释的原则\cite{stone1974,arlot2010,hastie2009}。
  \item \textbf{本模型保留了条件效应的证据路径。} 每一因素结论都可追溯至具体的匹配背景、共同温度与离散水平；同一因素在不同背景中反转时，模型报告反转而不强行汇总为一个全局主效应。局部二阶对比只在四格齐全时计算，并以删一温度检验符号是否由单点驱动。
  \item \textbf{本模型把优化与后续实验衔接起来。} 问题四将新增预算分配给高收益温区的空档、严格匹配结构的缺格和独立复验，既服务于当前推荐，也服务于下一轮可比较证据的形成；这与实验设计中以对照结构和重复性支撑解释的思想一致\cite{montgomery2017,box2005}。
\end{enumerate}

\subsection{模型不足}

\begin{enumerate}
  \item \textbf{本模型的优化域受实验覆盖限制。} 分段线性规则只在同一组合的相邻实测温度之间有效，候选因素级代理在留出高收益组合时误差较大。因此，未实验配方、跨组合连续设计空间和区间外温度均不能据此宣称最优。
  \item \textbf{本模型不能从无重复条件中分离纯实验误差。} 附件1大部分“组合--温度”单元只有一次观测；留一误差、删一温度敏感性和两次复验的样本方差分别刻画预测稳定性、单点驱动风险和基础重复性，均不等同于完整的显著性检验或总体方差估计。
  \item \textbf{本模型对可比较结构具有依赖性。} 严格匹配会主动舍弃无法控制其余因素的记录，以换取解释清晰度；当后续新增实验进入新的因素区域时，原有匹配块和局部四格未必仍存在，需要按新的实际设计重新枚举。
  \item \textbf{本模型不赋予机理因果解释。} 催化材料性质与反应路径在文献中具有复杂耦合\cite{pomalaza2016,hanspal2017,cimino2018}，而本文数据未包含表征、动力学或重复实验信息，故不能把条件差异解释为特定活性位或反应步骤的因果贡献。
\end{enumerate}

\subsection{模型改进方向}

\begin{enumerate}
  \item 在问题二已出现方向反转的匹配背景，以及问题三高收益候选附近，优先增加同条件重复实验；获得纯误差后，可将当前的条件差分扩展为带不确定性的区间比较。
  \item 在保持温度层和关键因素平衡的前提下补充正交或近正交设计点，再建立受层级约束的响应面或层次模型；模型拟合、交叉验证和优化应使用彼此隔离的验证单元。
  \item 对问题四第一阶段的新增实测结果，先更新观测排序、严格匹配块和决策裕度，再按预定规则确定复验对象；若出现新的高收益区域，应以实际新增点为起点重新设计后续温度加密，而不是沿用旧曲线的外推峰值。
\end{enumerate}
